\section{Nicht-funktionale Anforderungen} \todo{text}
% Zuverlässigkeit
% Wartbarkeit
% Modularität
% Erweiterbarkeit
Die nicht-funktionalen Anforderungen beschreiben qualitative Eigenschaften des Systems.

Das System soll eine hohe Zuverlässigkeit aufweisen. Die mittlere Betriebsdauer zwischen zwei Ausfällen soll mindestens 1000~h betragen.

Die Wartbarkeit soll durch eine strukturierte Softwarearchitektur sowie durch dokumentierte Schnittstellen unterstützt werden. Der Austausch einzelner Hardwarekomponenten soll ohne vollständige Demontage möglich sein.

Die Konstruktion soll modular aufgebaut sein. Anzeigeeinheit, Steuerplatine und Sensorik sollen getrennt ausgeführt werden. Erweiterungen, beispielsweise zusätzliche Informationsanzeigen, sollen ohne grundlegende Systemänderung integrierbar sein.

Die elektrische Sicherheit soll den geltenden Normen für Niederspannungsgeräte entsprechen. Die Erwärmung der Bauteile soll innerhalb zulässiger Grenzwerte bleiben.
